\documentclass[10pt,a4paper]{amsart}
\usepackage[margin=19mm]{geometry}
\usepackage{amsmath,amssymb,mathtools}
\usepackage[scr=rsfs,cal=euler]{mathalfa}
\usepackage{xfrac}
\usepackage[unicode,hidelinks,bookmarks=false]{hyperref}
\newcommand{\ie}{i.e.\ }
\newcommand{\cf}{cf.\ }
\newcommand{\resp}{respectively\ }
\newcommand{\Subsection}{\subsection}
\providecommand{\nobreakdash}{\nobreak\hbox{-}}
\setlength{\emergencystretch}{2em}
\multlinegap0pt
\usepackage{bm}
\usepackage{bbm}
\usepackage{dsfont}
\usepackage{xspace}
\usepackage{cancel}
\usepackage{mathtools}
\usepackage{amsthm}
\makeatletter
\@ifundefined{lemma}{\newtheorem{lemma}{Lemma}}{}
\makeatother
\title[Generalization Bounds for Transformer Channel Decoders]{Generalization Bounds for Transformer Channel Decoders}
\author{Qinshan Zhang, Bin Chen, Yong Jiang,, Shu-Tao Xia}
\date{Source: arXiv:2601.06969v1 (CC BY 4.0)}
\begin{document}
\maketitle

\noindent\emph{Source and licence.} Generalization Bounds for Transformer Channel Decoders. Version arXiv:2601.06969v1 (CC BY 4.0), distributed under the Creative Commons Attribution 4.0 International licence, as recorded in the arXiv licence metadata for this exact version and shown on the abstract page https://arxiv.org/abs/2601.06969v1. Full rights notice: licenses/arxiv-2601.06969-CC-BY-4.0.txt.\par\smallskip

\noindent\textit{Editorial scope.} Counted unit: the Lemma whose source label is \texttt{lemma: the sparsity of gradient df/dU} in the article (source0.tex lines 668--670, source label \texttt{lemma: the sparsity of gradient df/dU}), delivered together with the article's own complete original proof of it (source0.tex lines 672--722, one \texttt{proof} environment of 51 lines). No mathematics is written, repaired or re-derived: statement and proof are the article's own text line for line. Required context retained uncounted: none. Every definition, display and label of those lines is retained unchanged. Omitted from this package and not redistributed: 1--667, 671--671, 723--1432. Nothing inside a retained excerpt is omitted or altered: every retained body is delivered exactly as the article's lines are written, so no omission receipt of any kind accompanies this package. Citations: the retained excerpts carry no citation command at all, so no bibliography entry of the article is transcribed and no reference is redistributed. Reference adaptation: none was needed and none was made; every reference inside the delivered unit resolves inside it, so no unresolved reference or citation is produced. Printed slips kept literal: no printed word, symbol, hypothesis or number of the counted statement or of its proof was altered. Layout and class: the article's own class is \texttt{IEEEtran} with packages including amsmath, amsfonts, amssymb, amsthm, algorithmic, algorithm, array, textcomp, stfloats, url, verbatim, graphicx, cite, subfig; the wrapper is the lane's pinned \texttt{amsart} editorial wrapper at 10pt on A4, which restates the theorem-like environments the retained text needs and adds the article's own macro definitions that the retained text needs (none), with extra wrapper packages (bm, bbm, dsfont, xspace, cancel, mathtools). The delivered numbering is the unit's own. Assets: no retained span contains a figure environment, a graphics-inclusion command, a PDF-page inclusion, a table or a TikZ picture; no image file is copied into this package, the package declares no asset dependency and no image file is redistributed with it. Licence: version arXiv:2601.06969v1 of this article is distributed under the Creative Commons Attribution 4.0 International licence, as recorded in the arXiv licence metadata for this exact version and shown on the abstract page https://arxiv.org/abs/2601.06969v1. A full rights notice is delivered at licenses/arxiv-2601.06969-CC-BY-4.0.txt. Characters: every retained excerpt is pure ASCII, so the delivered engine, XeLaTeX with its default Latin Modern text font, reports no missing character.
\begin{lemma}\label{lemma: the sparsity of gradient df/dU}
    For the ECCT decoder function $f$, the gradient with respect to the pre-masking logits $\mathbf{U}$, denoted as $\nabla_{\mathbf{U}}f$, preserves the same sparsity pattern as $\Omega$.
\end{lemma}

\begin{proof}
$\nabla_{\mathbf{U}}f$ is computed via the chain rule:
\begin{equation}
    \begin{aligned}
    \frac{\partial f}{\partial U_{i j}}&=\sum_{k, l} \frac{\partial f}{\partial A_{k l}} \cdot \frac{\partial A_{k l}}{\partial U_{i j}}\\
    &=\sum_{k, l} \frac{\partial f}{\partial A_{k l}}\left(\sum_{m, n} \frac{\partial A_{k l}}{\partial S_{m n}} \cdot \frac{\partial S_{m n}}{\partial U_{i j}}\right).
    \end{aligned}
    \label{eq: df / dU_ij}
\end{equation}

The key point of our proof lies in analyzing the sparsity of the partial derivative $\frac{\partial A_{k l}}{\partial U_{i j}}$. Specifically, we aim to demonstrate that the Jacobian matrix inherits the sparsity structure of $\Omega$, ensuring that the gradient $\nabla_{\mathbf{U}}f$ remains $P$-sparse.

In \eqref{eq: df / dU_ij}, the term $\frac{\partial S_{m n}}{\partial U_{i j}}$ can be further written as :
\begin{equation}
    \frac{\partial S_{m n}}{\partial U_{i j}} = \delta_{mi} \delta_{nj}\mathbb{I}\left((i,j)\in \Omega \right),
    \label{eq: dSmn/dUij}
\end{equation}
where for any two non-negative integers $a$ and $b$, $\delta_{ab}$ denotes the Kronecker delta, which equals 1 if and only if $a=b$ and 0 otherwise. The indicator function $\mathbb{I}\left((i,j)\in \Omega \right)$ equals 1 if and only if the index $(i,j) \in \Omega$ and 0 otherwise. This implies that the partial derivative in \eqref{eq: dSmn/dUij} equals 1 if and only if $S_{m n}$ and $U_{i j}$ correspond to the same position (i.e., have the same index), which will not be masked, and equals 0 otherwise.

Therefore, $\frac{\partial A_{k l}}{\partial U_{i j}}$ in \eqref{eq: df / dU_ij} can be further simplified as follows:

\begin{equation}
    \frac{\partial A_{k l}}{\partial U_{i j}}=\frac{\partial A_{k l}}{\partial S_{i j}} \cdot \frac{\partial S_{i j}}{\partial U_{i j}}=\frac{\partial A_{k l}}{\partial S_{i j}} \cdot \mathbb{I}((i, j) \in \Omega).
\end{equation}

We now reduce the problem to analyzing $\frac{\partial A_{k l}}{\partial S_{i j}}$. Since $\mathbf{A}$ is obtained from $\mathbf{U}$ via the row-wise Softmax operation, the partial derivative $\frac{\partial A_{k l}}{\partial S_{i j}}$ can be expressed as:

\begin{equation}
    \frac{\partial A_{k l}}{\partial S_{i j}}=\delta_{k i} \cdot A_{k l}\left(\delta_{l j}-A_{i j}\right).
    \label{eq: dA_kl/S_ij}
\end{equation}

Clearly, when $k \neq i$, \eqref{eq: dA_kl/S_ij} equals 0. Therefore, it suffices to consider the case where $A_{k l}$ and $S_{i j}$ lie in the same row, which is consistent with the property of the row-wise Softmax function. Equation \eqref{eq: dA_kl/S_ij} can be further rewritten as:
\begin{equation}
    \frac{\partial A_{i l}}{\partial S_{i j}}= A_{i l}\left(\delta_{l j}-A_{i j}\right).
\end{equation}

Using the fact that $A_{i j}=0$ if $(i,j)\notin \Omega$, we can easily verify that $\frac{\partial A_{i l}}{\partial S_{i j}}=0$ if this position is masked. Therefore, \eqref{eq: df / dU_ij} can finally be written as:
\begin{equation}
    \begin{aligned}
        \frac{\partial f}{\partial U_{i j}} &= \sum_{l}\frac{\partial f}{\partial A_{il}} \cdot \frac{\partial A_{i l}}{\partial S_{i j}} \\
        &= \begin{cases}
            \sum_{l} \frac{\partial f}{\partial A_{il}}\cdot A_{i l}\left(\delta_{l j}-A_{i j}\right), & (i,j) \in \Omega \\
            0, & (i,j) \notin \Omega
        \end{cases}.
    \end{aligned}
\end{equation}

In summary, we have shown that the gradient $\nabla_{\mathbf{U}}f$ preserves the same sparsity pattern as $\Omega$. This completes the proof.

\end{proof}

\end{document}
